% !PRINTED 81
\weberpage{81}
%<<CONT

of \(a_\nu^{(\mu)}a_k^{(i)}\) in the expansion of \(A\), and \(A_\nu^{(\mu)}\) is the coefficient of \(a_\nu^{(\mu)}\); consequently \(A_{\nu,k}^{\mu,i}\) is the coefficient of \(a_k^{(i)}\) in the determinant \(A_\nu^{(\mu)}\).

According to this theorem, the so-called bordered determinant
\begin{equation}
        U=\begin{vmatrix}
        a_1^{(1)}&a_2^{(1)}&\cdots&a_n^{(1)}&u_1\\
        a_1^{(2)}&a_2^{(2)}&\cdots&a_n^{(2)}&u_2\\
        \vdots&\vdots&&\vdots&\vdots\\
        a_1^{(n)}&a_2^{(n)}&\cdots&a_n^{(n)}&u_n\\
        v_1&v_2&\cdots&v_n&q
        \end{vmatrix}
\tag{11}
\end{equation}
can be expanded with respect to the elements of the last row and column, and one obtains:
\begin{equation}
        U=qA-\sum^{i}\sum^{k}u_i v_k A_k^{(i)}.
\tag{12}
\end{equation}
This equation is obtained from (10) by replacing \(n\) by \(n+1\), and distinguishing the elements of the last row and column by a different notation.

For the higher subdeterminants also, notation by differential quotients is sometimes convenient, so that, for example,
\begin{equation}
        A_{\nu,k}^{i,\mu}=
        {\partial^2 A\over \partial a_\nu^{(i)}\,\partial a_k^{(\mu)}}
\tag{13}
\end{equation}
is put.

\begin{center}
{\large \S.\ 24.}

\vspace{0.5em}
{\large\gesperrt{Linear Homogeneous Equations.}}
\end{center}

The principal application of determinants, to which the whole theory owes its origin, is the solution of linear equations.

We shall attack the problem here at once in its most general form, since the special form affords scarcely any simplification and can afterwards easily be derived from the general result.

We consider a system of \(m\) equations of the first degree in which \(n\) unknowns \(x_1,x_2,\ldots,x_n\) occur homogeneously:
%CONT>>
